\chapter{the Loss Function}
\label{loss}


% commands
\newcommand{\eeg}{\mathbf{e}}
\newcommand{\att}{\mathbf{sa}}
\newcommand{\unatt}{\mathbf{su}}
\newcommand{\sima}[2]{\mathrm{sim}\bigl(\mathbf{\eeg}_{#1}, \mathbf{\att}_{#2}\bigr)}
\newcommand{\simu}[2]{\mathrm{sim}\bigl(\mathbf{\eeg}_{#1}, \mathbf{\unatt}_{#2}\bigr)}


% text
Embeddings are represented as: 
\begin{align*}
\eeg_i &\in \mathbb{R}^D \quad \text{(EEG embedding)}\\
\att_i &\in \mathbb{R}^D \quad \text{(attended speech embedding)}\\
\unatt_i &\in \mathbb{R}^D \quad \text{(unattended speech embedding)}
\end{align*}

\subsubsection*{Step 1: compute the similarity matrix}
To compare the embeddings in the latent space the measure of similarity is used. The similarity between two embeddings is defined as their inner product divided by a learnable temperature parameter $\tau$.

\[
\mathrm{sim}(\mathbf{\eeg_i},\mathbf{s_i}) = \frac{<\mathbf{\eeg_i},\mathbf{s_i}>}{\tau}
\]

When computing all possible combinations, two matrices are created. $S_A$ contains the matches between the eeg embeddings and the attended speech embeddings. $S_U$ contains eeg-unattended speech combinations.
\[
S_A = \begin{bmatrix}
\sima{0}{0} & \cdots & \sima{0}{N-1} \\
\vdots & \ddots & \vdots\\
\sima{N-1}{0} & \cdots & \sima{N-1}{N-1}
\end{bmatrix}
\]
\[
S_U = \begin{bmatrix}
\simu{0}{0} & \cdots & \simu{0}{N-1} \\
\vdots & \ddots & \vdots\\
\simu{N-1}{0} & \cdots & \simu{N-1}{N-1}
\end{bmatrix}
\]
\\
The diagonal of $S_A$ corresponds to correctly aligned EEG–attended speech pairs.
The diagonal of $S_U$ corresponds to aligned unattended speech, which is treated separately as a weak negative. Only $S_A$ is used in the CLIP loss computation. $S_U$ is excluded from the contrastive classification and is used solely in the weak-negative term. This means $S_U$ will in practice only need to be diagonal, as the rest of the combinations will not be used.

\subsubsection*{Step 2: compute the softmax probabilities}
For a given EEG embedding $\eeg_i$, the softmax probability $p_{e_i}$ measures how likely the attended speech $\att_i$ is relative to all other attended speech candidates in the batch. This formulation treats mismatched attended speech as contrastive negatives. 
\[
p_{e_i} =
\frac{\exp(S_A(i,i))}
{\sum_{j=1}^{B} \exp(S_A(i,j))}
\]
\\

For symmetry, we also compute the probability that $\eeg_i$ is the correct EEG match for $\att_i$:
\[
p_{sa_i} =
\frac{
\exp\big(\mathrm{sim}(\att_i,\eeg_i)/\tau\big)
}{
\sum_{j=1}^{B}\exp\big(\mathrm{sim}(\att_i,\eeg_j)/\tau\big)
}
\]

\subsubsection*{Step 3: Compute the the symmetric CLIP loss}
The loss is defined as:
\[
\mathcal{L}_{e} =
-\frac{1}{B}\sum_{i=1}^{B}\log p_{e_i}
\]

\[
\mathcal{L}_{sa} =
-\frac{1}{B}\sum_{i=1}^{B}\log p_{sa_i}
\]

\[
\mathcal{L}_{CLIP} = \frac{1}{2}\Bigl(
\mathcal{L}_{e}+
\mathcal{L}_{sa}
\Bigr)
\]

\subsubsection*{Step 4: Introducing the weak negative}
The temporally aligned unattended speech is treated as a weak negative using a margin ranking term. This margin-based term enforces that the EEG embedding is more similar to the attended speech than to the temporally aligned unattended speech by at least a margin $m$.
Unlike the CLIP loss, this term does not force unattended speech to be far from the EEG embedding, but only enforces a relative ordering, reflecting partial neural tracking of unattended speech.

\[
\mathcal{L}_{\mathrm{weak}} =
\frac{1}{B}\sum_{i=1}^{B}
\max\Bigl(0,\;
m - \bigl(\sima{i}{i}-\simu{i}{i}\bigr)
\Bigr).
\]

Finally, the total loss is:
\[
\mathcal{L}_{\mathrm{total}} =
\mathcal{L}_{CLIP}
+
\lambda \mathcal{L}_{\mathrm{weak}}
\]

In summary, the symmetric CLIP loss learns a shared embedding space by aligning EEG with attended speech using in-batch negatives, while the weak-negative term introduces task-specific structure by enforcing that attended speech is more similar to EEG than the temporally aligned unattended speech. This separation avoids treating unattended speech as a hard negative, which would contradict known neurophysiological evidence that EEG partially tracks both speech streams.

